\documentclass[a4paper,12pt]{article}

\usepackage[T2A]{fontenc}
\usepackage[utf8]{inputenc}
\usepackage[russian]{babel}

\usepackage{amsmath, amssymb}
\usepackage{geometry}
\usepackage{fancyhdr}
\usepackage{enumitem}

\geometry{margin=2cm}
\setlength{\parskip}{6pt}
\setlength{\parindent}{0pt}

\pagestyle{fancy}
\fancyhf{}
\rhead{11 класс}
\cfoot{\thepage}

\begin{document}

\begin{center}
    {\Large \textbf{Целые числа}}\\[0.5em]
\end{center}
    
\begin{enumerate}[label=\textbf{Задача \arabic*.}]
\item Найдите все такие пары целых чисел m и $n > 2$, что $((n-1)! - n) \cdot (n-2)! = m(m - 2)$.
\textbf{Решение:} Заметим, что
\begin{align*}
((n-1)!-1)((n-2)!-1) &= (n-1)!(n-2)!-(n-1)!-(n-2)!+1 \\
&= ((n-1)!-n)(n-2)!+1 = m^2-2m+1 = (m-1)^2.
\end{align*}
Пусть $n > 4$. Заметим, что числа $(n-1)!-1$ и $(n-2)!-1$ взаимно просты. Предположим, что это не так, и оба этих числа делятся на простое число $p$. Тогда число $(n-1)! - 1 - ((n-2)! - 1) \cdot (n-1) = n-2$ тоже делится на $p$. Тогда $(n-2)!$ делится на $p$, а $(n-2)!-1$ не кратно $p$, противоречие.

Таким образом, произведение взаимно простых чисел $(n-1)!-1$ и $(n-2)!-1$ --- точный квадрат, тогда и каждое из них точный квадрат. Однако, число $(n-1)!-1$ при $n > 4$ даёт остаток 3 при делении на 4, поэтому оно точным квадратом быть не может.

Остаётся разобрать случаи $n \le 4$. При $n=4$ получается $(m-1)^2=5$, решений нет. При $n=3$ мы получаем: $(m-1)^2 = 0$, что даёт единственное решение $m=1, n=3$.

\item Числа от 1 до 999999 разбиты на две группы: в первую отнесено каждое число, для которого ближайшим к нему квадратом является квадрат нечетного числа, во вторую — числа, для которых ближайшими являются квадраты четных чисел. В какой из групп сумма чисел
больше? \\
\textbf{Решение:} Эти суммы одинаковы.

Разобьем числа от $n^2$ до $(n+1)^2 - 1$ на две группы:
\begin{align*}
A_n &= \{n^2, n^2+1, \dots, n^2+n\} \\
B_n &= \{n^2+n+1, n^2+n+2, \dots, n^2+2n\}
\end{align*}

Для чисел группы $A_n$ ближайшим квадратом является $n^2$, а для группы $B_n$ ближайшим является $(n+1)^2$ --- квадрат другой четности. Обозначим суммы чисел в этих группах как $S(A_n)$ и $S(B_n)$ соответственно.

Рассмотрим разность сумм:
\begin{align*}
S(B_n) - S(A_n) &= [(n^2+n+1) - (n^2+1)] \\
&\quad + [(n^2+n+2) - (n^2+2)] \\
&\quad + \dots \\
&\quad + [(n^2+2n) - (n^2+n)] - n^2 \\
&= \underbrace{n+n+\dots+n}_{n \text{ раз}} - n^2 \\
&= n \cdot n - n^2 = 0
\end{align*}

Таким образом, $S(A_n) = S(B_n)$.

Заметим, что все множество чисел от 1 до 999999 разбивается на непересекающиеся пары $(A_1, B_1), (A_2, B_2), \dots, (A_{999}, B_{999})$. Поскольку в каждой паре суммы чисел равны, то и общие суммы для чисел с ближайшими четными и нечетными квадратами также равны.
\item Докажите, что уравнение $m + s! = p$, где все три числа целые, m - составное и p - простое, имеет бесконечно много решений при фиксированном s. \\
\textbf{Решение:} Для любого натурального k существует k подряд идущих составных чисел (достаточно рассмотреть последовательность чисел $(k+1)! + 2, \, (k+1)! + 3, \, \ldots \, , (k+1)! + (k+1)$). Рассмотрим первое просто число после последовательных $s!$ составных чисел - оно и является одним из решений уравнения. Для получения бесконечного числа решений достаточно рассмотреть $s!, \, 2\cdot s!, \, 3\cdot s!,\, \ldots \, ,\ell \cdot s!, \, \ldots$ последовательных составных чисел, брать следующее после них простое число и получать новое решение. 

\item Назовём главными делителями составного числа n два наибольших его натуральных делителя, отличных от n. Составные натуральные числа a и b таковы, что главные делители числа a совпадают с главными делителями числа b. Докажите, что $a = b$. \\
\textbf{Решение:} Пусть $n > k$ --- главные делители числа $a$; тогда $\frac{a}{n}$ и $\frac{a}{k}$ --- два наименьших делителя числа $a$, больших единицы. Пусть $p$ --- наименьший простой делитель числа $a$, а $q$ --- наименьший простой делитель $a$, кроме $p$ (если такой существует). Тогда $\frac{a}{n} = p$. Далее, $\frac{a}{k}$ --- либо простое число (тогда это $q$), либо составное (значит p - единственный простой делитель числа $a$). Во втором случае единственным простым делителем числа $\frac{a}{k}$ является $p$ (иначе у нас есть новый простой делитель $a$), и потому $\frac{a}{k} = p^2$. Тогда, если q существует, то главные делители - $\frac{a}{p}$ и $\frac{a}{q}$, а если простого q не существует, то $\frac{a}{p}$ и $\frac{a}{p^2}$.

Итак, главные делители числа $a$ --- это либо $\frac{a}{p}$ и $\frac{a}{q}$, либо $\frac{a}{p}$ и $\frac{a}{p^2}$. Покажем теперь, что по двум главным делителям $n > k$ составное число $a$ восстанавливается однозначно (откуда и следует требуемое). Если $n$ кратно $k$, то выполнен второй случай, и тогда $a = n^2/k$. Иначе выполнен первый случай, и тогда $a = \text{НОК}(n,k)$ (ибо \text{НОД}(n,k) = $\frac{a}{pq}$, это легко доказать: если $\frac{a}{p}$ = $mq$, то есть \text{НОД}(n,k) = \text{НОД}($mq$, $mp$) = $m = \frac{a}{qp}$).

\item Найдите все пары чисел $(p, n)$, где $p$--- простое, а $n$--- целое и при этом $p^4 + 211 = n^2 + 5pn$. В ответе укажите значения $n$. Если их несколько, перечислите их в любом порядке через запятую. \\
\textbf{Решение:} Согласно малой теореме Ферма, если $p \neq 5$, число $p^4$ даёт остаток 1 при делении на 5. Число 211 также даёт остаток 1 при делении на 5. И $5pn$ делится на 5. Значит, если $p \neq 5$, мы получаем, что $n^2$ даёт остаток 2 при делении на 5, что невозможно.

Осталось разобрать случай $p = 5$. В этом случае нам надо решить квадратное уравнение
\[ n^2 + 25n - 625 - 211 = 0 \]
\[ n_{1,2} = \frac{-25 \pm \sqrt{3125 + 4 \cdot 211}}{2} \]
Откуда получаем два решения: $(p, n) = (5, -44)$ и $(p, n) = (5, 19)$.

\item Найдите все тройки целых чисел $x, y, z$ такие, что $5x^2 + 9y^2 = 13z^2$. \\
\textbf{Решение:} Для начала заметим, что подходит очевидная тройка чисел $(0, 0, 0)$. Если $x = 0$, то получившееся равенство не имеет решений (кроме тривиального), потому что степени вхождения 13 в обе части не будут совпадать. Аналогично, если какая-то другая одна переменная равна 0. Поэтому в дальнейшем рассматриваем ненулевые числа. Теперь давайте доказывать, что других троек не существует. Предположим противное. Тогда рассмотрим тройку $(x, y, z)$ такую, что степень вхождения 5 в $x$ минимальная возможная.
Рассмотрим уравнение по модулю 5. Имеем
\[ 9y^2 \equiv 13z^2 \pmod{5} \]
Квадраты целых чисел могут давать остатки 0, 1 и $-1$ по рассматриваемому модулю. Предположим, что никакое из чисел $y, z$ не кратно 5. Но тогда
\[ 1 \equiv |y^2| \equiv |9y^2| \equiv |13z^2| \equiv |2z^2| \equiv 2 \pmod{5} \]
что невозможно. Таким образом, по крайней мере одно из чисел $y, z$ кратно 5, но тогда 5 делит каждое из них.
Наконец, $5x^2 = 13z^2 - 9y^2$ кратно 25, следовательно, $x$ кратно 5. Но тогда тройка $(\frac{x}{5}, \frac{y}{5}, \frac{z}{5})$ так же является решением. Таким образом, мы получили противоречие.
\item Решить уравнение в целых числах: $m^2 + k^2 = 2024^n + 33$. \\
\textbf{Решение:} Пусть $n \ge 2$. Тогда $2024^n + 33$ делится на 11, поскольку каждое из чисел 2024 и 33 кратно 11, но не делится на $11^2$, т.к. первое слагаемое кратно $11^2$, а второе --- нет.

Пусть число $x$ дает остаток 0, 1, 2, 3, 4, 5, 6, 7, 8, 9, 10 при делении на 11, тогда число $x^2$ дает соответственно остаток 0, 1, 4, 9, 5, 3, 5, 9, 4, 1 при делении на 11. Докажем, что если хотя бы одно из чисел $m$ и $k$ не делится на 11, то и число $m^2 + k^2$ не делится на 11. Предположим обратное, тогда сумма остатков чисел $m^2$ и $k^2$ равна 11, следовательно, ровно одно из чисел $m^2$ и $k^2$ даёт четный остаток при делении 11, а значит, соответствующий квадрат даёт остаток 0 или 4 при делении на 11, но тогда второй остаток равен 0 или 7, что невозможно. Таким образом, каждое из чисел $m$ и $k$ кратно 11, следовательно, каждое из чисел $m^2$ и $k^2$ кратно $11^2$, таким образом, $m^2+k^2$ кратно $11^2$, но $2024^n + 33$ не кратно $11^2$.

Тогда $n=0$ или $n=1$.

Пусть $n = 1$. Тогда $2024^n + 33 = 2024 + 33 = 2057 = 11^2 \cdot 7$, следовательно, $m^2 + k^2$ кратно $11^2$, а значит, как мы показали выше, каждое из чисел $m$ и $k$ кратно 11. Пусть $m = 11a$, $k = 11b$, где $a$ и $b$ являются целыми числами, следовательно, $a^2 + b^2 = 17$. Легко убедиться, что всеми решениями $(a,b)$ данного уравнения являются неупорядоченные пары $(\pm 1, \pm 4)$. Следовательно, все пары решений $(m,k)$ это $(\pm 11, \pm 44), (\pm 44, \pm 11)$.

Пусть $n=0$. Тогда $2024^n + 33 = 1 + 33 = 34$. Если каждое из чисел $m$ и $k$ не превосходит по модулю 4, то сумма их квадратов не превосходит 32, следовательно, наибольшее из чисел $m$ и $k$ по модулю не меньше 5. С другой стороны, если какое-то из чисел по модулю больше 5, то его квадрат не меньше 36, что невозможно. Таким образом, в паре чисел $(m,k)$ хотя бы одно равно 5 по модулю, тогда второе равно 3 по модулю. Тем самым, мы показали, что все пары решений $(m,k)$ есть $(\pm 5, \pm 3), (\pm 3, \pm 5)$.

\item Найдите все многочлены $P(x)$ над $\mathbb{Z}$ такие, что $P(P(x)+x)$ является простым числом при бесконечном количестве целых x. \\
\textbf{Решение:} Заметим, что $P(P(x) + x) - P(x)$ делится на $P(x) + x - x = P(x)$ при каждом натуральном $x$. При этом раз в бесконечном количестве натуральных точек $P(P(x) + x)$ --- простое, то $P(x)$ в бесконечном количестве натуральных точек или равен 1, или равен $-1$, или равен $P(P(x) + x)$, или $-P(P(x) + x)$. Первые два случая нам не подходят.
Предположим, что $P(x)$ тождественно равен $P(P(x) + x)$. Если степень многочлена $P(x)$ больше 1, то степень $P(P(x)+x)$ больше степени $P(x)$, поэтому $P(x) = c$, или $P(x) = ax + b$. В первом случае нам подходят все простые $c$. Во втором случае имеем $ax + b = a(a+1)x + b(a+1)$, откуда $a \cdot (ax+b) = 0$, что невозможно при $a \neq 0$.
Если же $P(x)$ тождественно равен $-P(P(x)+x)$, то достаточно рассмотреть случай $P(x) = ax+b$ ($a \neq 0$). Тогда $ax+b = -a(a+1)x - b(a+1)$, откуда $(a+2)(ax+b)=0$. То есть $a=-2$. Осталось лишь понять, что только при нечётных $b$ выражение $-2x+b$ является простым в бесконечном количестве целых точек.

\item Натуральное число $N$ представляется в виде $N = a_1 - a_2 = b_1 - b_2 = c_1 - c_2 = d_1 - d_2$, где $a_1$ и $a_2$ --- квадраты, $b_1$ и $b_2$ --- кубы, $c_1$ и $c_2$ --- пятые степени, а $d_1$ и $d_2$ --- седьмые степени натуральных чисел. Обязательно ли среди чисел $a_1$, $b_1$, $c_1$ и $d_1$ найдутся два равных? \\
\textbf{Решение:} Приведём пример числа $N$, для которого все указанные числа будут различными. Положим
\[ N = (3^2-2^2)^{105}(3^3-2^3)^{70}(3^5-2^5)^{126}(3^7-2^7)^{120}. \]
Тогда
\[ N = M_2^2(3^2-2^2) = M_3^3(3^3-2^3) = M_5^5(3^5-2^5) = M_7^7(3^7-2^7) \]
при некоторых натуральных $M_2, M_3, M_5$ и $M_7$, не делящихся ни на 2, ни на 3. Отсюда
\begin{align*}
N &= (3M_2)^2 - (2M_2)^2 = (3M_3)^3 - (2M_3)^3 = \\
  &= (3M_5)^5 - (2M_5)^5 = (3M_7)^7 - (2M_7)^7.
\end{align*}
Даже все восемь чисел, участвующих в представлениях, различны, поскольку у любых двух из них разная степень вхождения либо двойки, либо тройки.

\item На свадьбу пришли несколько супружеских пар, у каждой из которых было от 1 до 10 детей. Тамада выбирал одного ребёнка, одну маму и одного папу из трёх разных семей и заставлял их участвовать в странных конкурсах. Оказалось, что у него было ровно 3630 способов выбрать нужную тройку людей. Сколько всего могло быть детей на этом прекрасном вечере? \\
\textbf{Решение.} Пусть на свадьбе было $p$ супружеских пар и $d$ детей (из условия, $d \le 10p$). Тогда каждый ребёнок состоял в $(p-1)(p-2)$ тройках: маму можно было выбрать из одной из $p-1$ супружеских пар, а при зафиксированном выборе мамы папу можно было выбрать из одной из $p-2$ оставшихся пар. Значит, общее количество троек равно $d \cdot (p-1)(p-2) = 3630$.
Поскольку $d \le 10p$, получаем $3630 \le 10p^3$, то есть $p^3 \ge 363 > 7^3$. Значит, $p \ge 8$.
Далее, число $3630 = 2 \cdot 3 \cdot 5 \cdot 11^2$ имеет два делителя $p-1$ и $p-2$, отличающиеся на 1. Если один из этих делителей делится на 11, то другой даёт остаток 1 или 10 при делении на 11. Тогда он взаимно прост с 11, а значит, он делит $2 \cdot 3 \cdot 5 = 30$ и при этом не меньше 10. Нетрудно видеть, что этим делителем может быть только 10; тогда $p-2 = 10, p-1 = 11$ и $d = 3630/110 = 33$.
Если же оба числа $p-2$ и $p-1$ не делятся на 11, то число $2 \cdot 3 \cdot 5 = 30$ делится на их произведение; это противоречит тому, что $p \ge 8$.

\item Коэффициенты $a, b, c$ квадратного трёхчлена $f(x) = ax^2 + bx + c$ --- натуральные числа, сумма которых равна 2000. Паша может изменить любой коэффициент на 1, заплатив 1 рубль. Докажите, что он может получить квадратный трёхчлен, имеющий хотя бы один целый корень, заплатив не более 1050 рублей.

\textbf{Решение:} Покажем, что можно из исходного трёхчлена получить новый, имеющий хотя бы один целый корень, изменив коэффициенты $a, b, c$ суммарно даже не более, чем на 1022.

Если коэффициент $c$ не больше 1022, то, сделав его равным нулю, мы получим искомый трёхчлен $f_1(x) = ax^2+bx$, имеющий целый корень $x=0$.

Пусть теперь $c \ge 1023$; тогда из условия $a+b \le 977$. Рассмотрим два последовательных квадрата, между которыми находится $c$: $m^2 \le c < (m+1)^2$. Поскольку $c < 2000 < 45^2$, имеем $m \le 44$. Тогда одна из разностей $c-m^2$ и $(m+1)^2-c$ не превосходит $\frac{1}{2}((m+1)^2 - m^2) = \frac{1}{2}(2m+1)$, то есть она не больше $m \le 44$. Итак, найдётся натуральное $k$, для которого $|c-k^2| \le 44$. Заменив теперь $a$ на $-1$, $b$ на $0$, а $c$ на $k^2$, мы изменим коэффициенты суммарно не более, чем на $(a+b+1) + |c-k^2| \le 978 + 44 = 1022$, и получим трёхчлен $f_2(x) = -x^2+k^2$, имеющий целый корень $x=k$.

\item Существует ли $f: \mathbb{Z} \longrightarrow \mathbb{Z}$ такая, что $f((x^2+1)\, \text{ mod }\, 7) = (f(x)^2 + 1) \text{ mod 11}$? \\
\textbf{Решение:} Стандартным ходом при решении задач на функциональные уравнения является подставить какое-то значение переменной, при котором два часто возникающих и не равных друг-другу тождественно выражения оказываются равны, и посмотреть, какие следствия из этого удастся вывести. Применительно к данной задаче на роль такой подстановки просится значение $x_0$, для которого выполнялось бы $x_0 = x_0^2 + 1 \pmod{7}$.

Задумаемся, а существует ли такое $x_0$? Условие равносильно квадратному уравнению в остатках: $x_0^2 - x_0 + 1 \equiv 0$ (в этом абзаце все сравнимости по модулю 7), эквивалентно $x_0 \equiv \frac{1 \pm \sqrt{-3}}{2} = \frac{1 \pm \sqrt{-3+7}}{2} = \{\frac{3}{2}, \frac{-1}{2}\} = \{\frac{3}{2}, \frac{-1+7}{2}\} = \{5, 3\}$. Или можно было просто перебором остатков, благо их всего 7, убедиться, что любой из 3 и 5 подходят.

Что же нам дает равенство $3 = 3^2 + 1 \pmod{7}$? Просится от обоих частей взять функцию $f$, а затем воспользоваться условием задачи. Имеем: $f(3) = f((3^2+1) \pmod{7}) = (f(3)^2+1) \pmod{11}$. Чтобы подчеркнуть полученное, обозначим $f(3) = y$ и выбросим среднюю часть: $y = (y^2 + 1) \pmod{11}$. Отсюда следует $y^2 - y + 1 \equiv 0$ (в этом абзаце все сравнимости по модулю 11), отметим что это именно следствие а не равносильность. Выясним, имеет ли сравнимость решения, действуя стандартно $y \equiv \frac{1 \pm \sqrt{-3}}{2}$. А извлекается ли квадратиый корень из -3 по модулю 11? Заметим что $1^2 \equiv (-1)^2 \equiv 1$, $2^2 \equiv (-2)^2 \equiv 4$, $3^2 \equiv (-3)^2 \equiv 9$, $4^2 \equiv (-4)^2 \equiv 5$ и $5^2 \equiv (-5)^2 \equiv 3$. Мы перебрали все остатки, среди квадратов не нашлось -3, значит корень не извлекается, значит уравнение $y^2 - y + 1 \equiv 0$ не имеет решений.

Итак, требуемой функции $f$ не существует.

\item Для таблички $n\times n$ рассматриваем семейство квадратов $2 \times 2$, состоящих из клеток таблицы, и обладающее свойством: для любого квадрата семейства найдется покрытая им клетка, непокрытая никаким другим квадратом из семейства. Через $f(n)$ обозначим максимальное количество квадратов в таком семействе. Для какого наименьшего C неравенство $f(n)\leq C n^2$ верно при любом n? \\
\textbf{Решение:} Во-первых, докажем что $f(n) \le \frac{1}{2}n^2$. Для этого полезно доказать более сильное утверждение: для произвольной фигуры из $S$ клеток количество квадратов $2 \times 2$ в семействе, таком что все квадраты лежат в фигуре и для любого квадрата найдется клетка, покрытая только им, не превосходит $S/2$. Рассмотрим два случая: для семейства найдется клетка A, покрытая четырьмя квадратами, и случай, когда такой клетки не найдется.

Если такая клетка A нашлась, то рассмотрим четыре покрывающих ее квадрата. Они образуют квадрат $3 \times 3$ с клеткой A в центре. И поскольку в каждом из четырех квадратов $2 \times 2$ должна быть клетка, покрытая только им --- это четыре угловые клетки квадрата $3 \times 3$, поскольку все остальные покрыты хотя бы дважды. Но тогда никакой другой квадрат $2 \times 2$ из семейства не покрывает клетки квадрат $3 \times 3$, иначе он покрывает и угловую клетку, а она должна быть покрыта только один раз. Итак, все остальные квадраты лежат в множестве площади $S - 9$. В этот момент доказательства еще не поздно решить, что на самом-то деле мы ведем индукцию по $S' :)$, благо база тривиальна. Итак, всего квадратов в семействе оказывается не больше $4 + \frac{S-9}{2} = \frac{S-1}{2} < S/2$.

Пусть клетки, покрытой четырьмя квадратами из семейства, не найдется. Поместим в каждую клетку множества единичный заряд. Теперь пусть каждая клетка, покрытая $k$ квадратами из семейства, отдает каждому из этих квадратов по $1/k$ заряда (таким образом, раздаст весь свой заряд). Тогда каждый квадрат семейства получил заряд не меньше 2, потому что минимум от одной клетки получил 1, и от остальных получал не меньше 1/3 от каждой. Итого, всего полученного заряда не меньше чем дважды число квадратов в семействе, а отданного заряда не больше $S$, итого квадратов в семействе не больше $S/2$.

Теперь построим пример, доказывающий, что $f(n) \ge \frac{1}{2}n^2 - 4n$, следовательно неравенство $f(n) \le Cn^2$ при $C < \frac{1}{2}$ неверно при всех достаточно больших $n$.

Возьмем бесконечную клетчатую плоскость и покрасим ее в два цвета следующим образом: выберем одно из двух направлений диагонали, покрасим все клетки каждой диагонали в один цвет: две диагонали в белый, следующие две в черный и так далее с периодом 4. Теперь выберем квадрат $n \times n$, в который черных клеток попало не меньше чем белых, то есть хотя бы $n^2/2$. Теперь на каждую черную клетку внутри квадрата $n \times n$ положим квадрат $2 \times 2$ так, чтобы кроме этой черной клетки квадрат содержал только белые (это можно сделать единственным образом). Удалим все квадраты $2 \times 2$, частично вылезшие за границы квадрата $n \times n$, их не больше $4n$. Требуемое семейство построено.

\item Из натурального числа $n$ разрешается получить либо число $n^2 + 2n$, либо число $n^3 + 3n^2 + 3n$. Два натуральных числа называются совместимыми, если из них можно получить одно и то же число с помощью некоторого количества таких операций. Найдите все числа, совместимые с числом 2018.

\textbf{Решение:} Сделаем замену $k = n + 1$ и будем считать, что мы преобразуем число $k$, которое может принимать значения натуральных чисел, кроме единицы. Замена $n \to n^2 + 2n$ для $k$ соответствует замене $f_1: k = n + 1 \to n^2 + 2n + 1 = k^2$. Вторая замена соответствует $f_2: k \to k^3$. Заметим, что для любого $k$ верно $f_1(f_2(k)) = f_2(f_1(k))$. Таким образом, если мы применяем несколько раз операции $f_1$ и $f_2$ к числу $k$, неважен порядок, а важно только количество операций.

Допустим, числа $k_1$ и $k_2$ эквивалентны. Тогда применением операций к одному и другому числам несколько раз можно получить одно и то же число, то есть $k_1^{2^{l_1}3^{m_1}} = k_2^{2^{l_2}3^{m_2}}$. Таким образом, все натуральные числа, эквивалентные заданному $k$, имеют вид $k^{2^{q_1}3^{q_2}}$ для рациональных $q_1, q_2$. Соответственно, для $n=2018$ все совместимые с ним числа будут иметь вид $2019^{2^{q_1}3^{q_2}} - 1$. Число $2019=3 \cdot 673$ не является степенью натурального числа выше первой. Таким образом, рациональные числа $q_1$ и $q_2$ должны быть целыми, для любых целых $q_1$ и $q_2$ мы получаем совместимые с 2018.

Ответ - числа вида $2019^{2^n \cdot 3^k} - 1$ для неотрицательных целых $k$ и $n$.

\item Натуральное число $n$ таково, что число $2n-1$ является простым. Докажите, что среди любых $n$ различных чисел натуральных $a_1, \ldots, a_n$ найдутся два различных числа $a_i$ и $a_j$ с условием 
$$
\dfrac {(a_i+a_j)}{\text{НОД}(a_i, a_j)} \geq 2n-1.
$$


\item Таблица $2 \times 2024$ заполнена целыми числами, причём в первой строке стоят числа из набора $\{1, ..., 2023\}$. Оказалось, что какие бы два столбца мы ни выбрали, разность их чисел из первой строки делится на разность их чисел из второй строки. Известно, что все числа во второй строке попарно различны. Обязательно ли тогда все числа в первой строке равны между собой? \\
\textbf{Решение:}  \\
\textbf{Лемма.} Пусть $k$ и $l$ -- натуральные числа, причём $\text{НОК}(k,l) < k+l$. Тогда $k+l$ делится на $k$ или на $l$.

\textbf{Доказательство.} Пусть $d = \text{НОД}(k,l)$. По условию $kl < d(k+l)$, т.е. $(k-d)(l-d) < d^2$. Один из множителей меньше $d$ и делится на $d$, т.е. равен нулю. А $k+l$ делится на $d$. $\square$

Условие делимости не изменится, если все числа строки уменьшить на одно и то же целое число. Поэтому можно считать, что верхние числа находятся в пределах от 0 до 2022, а нижние от -- 0 до $M$ (причём и 0, и $M$ присутствуют). Ясно, что $M > 2022$. Разность чисел над 0 и $M$ делится на $M$, поэтому эти числа равны (пусть это число $b$).

Предположим, есть столбец вида $\begin{pmatrix} b+d \\ k \end{pmatrix}$, где $d \neq 0$. Ясно, что $|d| \le 2022$. Сравнивая этот столбец со столбцами $\begin{pmatrix} b \\ 0 \end{pmatrix}$ и $\begin{pmatrix} b \\ M \end{pmatrix}$, видим, что $d$ делится как на $k$, так и на $M-k$.
Поэтому $d$ делится на $\text{НОК}(k, M-k)$; по лемме $M$ делится на $k$ или на $M-k$. Как известно, у числа $M$ не более $2\sqrt{M}$ делителей, столько же дополнений этих делителей до $M$, значит, всего в верхней строке не более $4\sqrt{M}$ чисел, отличных от $b$. С другой стороны, и $k$, и $M-k$ не больше 2022, поэтому $M \le 4044$. Итак, в верхней строке не более $4\sqrt{4044} < 300$ чисел, отличных от $b$. Зафиксируем один столбец вида $\begin{pmatrix} b+d \\ k \end{pmatrix}$ и рассмотрим произвольный столбец вида $\begin{pmatrix} b \\ k+q \end{pmatrix}$. По условию $d$ делится на $q$. Значит, в первой строке не более $4\sqrt{d} \le 4\sqrt{2022} < 180$ чисел $b$ ($q$ может быть как положительным, так и отрицательным). Но $300 + 180 < 2024$. Противоречие.
\end{enumerate}
\end{document}